\documentstyle[%


righttag%


,11pt%


,amssymb%


,russian%               remove in Eng.


,izvru%                 Set izven% in Eng.


% ,izvru-ks             for brief communications


]{amsart}





%       Math definitions





\newcommand{\col}{\operatorname{col}}


\newcommand{\spc}[1]{\bold {#1}}


\newcommand{\eqdef}{\overset{\mathrm{def}}{=}{}}


\newcommand{\vraisup}{\operatornamewithlimits{vrai\,sup}}


\newcommand{\vraiinf}{\operatornamewithlimits{vrai\,inf}}


\newcommand{\tts}[1]{}





%\hoffset=-15mm%                  For Eng.


\setcounter{page}{3}


\god{2000}


\nomer{4~(455)} %                        Remove in Eng.


%\nomer{Vol.\,44, No.\,4}%               Set in Eng.


%\str{\pageref{begin}--\pageref{end}}%  Set in Eng.


\udk{517.929}





\author{\it Н.В.\,Азбелев, П.М.\,Симонов}


% NB:   ^^ remove in Eng.


\title{Устойчивость уравнений\\ с запаздывающим аргументом. II}


\thanks{Работа выполнена при финансовой поддержке


Российского фонда фундаментальных исследований


(96-15-96195, 99-01-01278), International Soros Science Education Program


(EP-55, d99-1245) и Конкурсного


центра по исследованиям в области фундаментального естествознания,


Санкт-Петербург.}





\date{}


%\pagestyle{myheadings}%         Set in Eng.


\pagestyle{plain} %              Remove in Eng.


\begin{document}


%\thispagestyle{plain}%          Set in Eng.


\maketitle


\label{begin}





В предыдущих работах \cite{1}, \cite{2} на основе идей теории


``абстрактного функционально-диф\-фе\-рен\-ци\-аль\-ного


уравнения'' \cite{3}--\cite{5} была предложена новая концепция


устойчивости для линейных обыкновенных дифференциальных уравнений и


их обобщений. Ниже упомянутая концепция и результаты работ


\cite{6}--\cite{9} развиваются для уравнений квазилинейных.





Обозначим через $\spc{C}$ банахово пространство непрерывных и


ограниченных


функций $x:[0,\infty)\to\spc{R}^n$ с нормой


$\|x\|_{\bold C}\eqdef\sup\limits_{t\ge 0}|x(t)|$, где $|\cdot|$ ---


норма в ${\spc{R}}^n$.


Пусть, далее, пространство ${\spc{C}}_0$ и {\it весовое} пространство


${\spc{C}}_\gamma$ при $\gamma>0$ определяются равенствами


$$


{\spc{C}}_0\eqdef\{x\in{\spc{C}}:\lim_{t\to\infty}|x(t)|=0,


\ \|x\|_{{\bold C}_0}\eqdef \|x\|_{\bold C}\},\quad


{\spc{C}}_\gamma\eqdef\{x\in{\spc{C}}:\|x\|_{{\bold{C}}_\gamma}\eqdef


\sup_{t\ge 0}e^{\gamma t}|x(t)|<\infty\}.


$$


Пространства ${\spc{C}}_0$ и


${\spc{C}}_\gamma$ банаховы.





Задача Коши


\begin{equation}


\Dot x(t)=f(t,x(t)),\;x(0)=\alpha,\quad t\ge 0,


\end{equation}


эквивалентна уравнению


\begin{equation}


x(t)=\int_0^tf(s,x(s))\,ds+\alpha.


\end{equation}


Поэтому приведенные в (\cite{10}, c.\,9) определения устойчивости


уравнения $\Dot x(t)=f(t,x(t))$ в случае $f(t,0)\equiv 0$


можно переформулировать следующим образом: тривиальное решение задачи


(1) устойчиво по Ляпунову (асимптотически, экспоненциально), если


уравнение (2) имеет в пространстве ${\spc{C}}$ (в ${\spc{C}}_0$,


${\spc{C}}_\gamma$ соответственно) единственное решение при каждом


$\alpha\in{\spc{R}}^n$ и это решение непрерывно зависит от $\alpha$ в


метрике этого пространства. Другими словами, данный тип устойчивости


--- это корректная разрешимость уравнения (2) в данном пространстве.


В \cite{10} рассмотрена также ``устойчивость


относительно постоянно действующих возмущений $\eta$'' ---


непрерывная зависимость решения задачи


$$


\Dot x(t)=f(t,x(t))+\eta(t),\;x(0)=\alpha,\quad t\ge 0,


$$


от возмущений $\eta$.





Рассмотрим теперь общее функционально-дифференциальное уравнение


${\cal L}x={\cal F}x$ в предположении, что ${\cal


L}:{\spc{D}}\to\spc{L}$ --- линейный, а ${\cal


F}:{\spc{D}}\to\spc{L}$ --- нелинейный вольтерровы операторы \cite{3},


\cite{4}. Здесь $\spc{L}$ --- линейное пространство классов эквивалентных


локально суммируемых функций $z:[0,\infty)\to\spc{R}^n$, $\spc{D}$


--- линейное пространство всех абсолютно непрерывных на каждом


конечном отрезке функций $x:[0,\infty)\to\spc{R}^n$. Пусть, далее,


задано банахово пространство $\spc{B}$ с нормой


$\|\cdot\|_{\bold{B}}$, элементы которого принадлежат пространству


$\spc{L}$, а решение $x\in{\spc{D}}$ задачи Коши ${\cal


L}_0x=z$, $x(0)=\alpha$ для ``модельного'' уравнения ${\cal L}_0x=z$ с


линейным вольтерровым оператором ${\cal L}_0:{\spc{D}}\to\spc{L}$


при любом $z\in\spc{L}$ имеет


представление в виде формулы Коши


$x={\cal W}z+{\cal U}\alpha$ \cite{3}, \cite{4}. Здесь вольтерров


оператор ${\cal W}:\spc{L}\to\spc{D}$ --- оператор Коши модельного


уравнения, а ${\cal U}:\spc{R}^n\to\spc{D}$ --- оператор,


порожденный фундаментальной матрицей $U$ решений однородного


уравнения ${\cal L}_0x=0$: $({\cal U}\alpha)(t)\eqdef U(t)\alpha$.


Банахово пространство ${\spc{D}}_0$ определим равенствами


$$


\spc{D}_0\eqdef{\cal W}\spc{B}\oplus{\cal U}{\spc{R}}^n,\quad


\|x\|_{\bold{D}_0}\eqdef\|{\cal L}_0x\|_{\bold B}+|x(0)|.


$$





Развивая идею понятия устойчивости как корректной разрешимости


задачи Коши в данном пространстве, введем





\begin{defn}


Будем говорить, что уравнение ${\cal L}x={\cal F}x$\;


{\it ${\spc{D}}_0$-устойчиво}


(обладает {\it ${\spc{D}}_0$-свойством}),


если задача Коши


\begin{equation}


{\cal L}x={\cal F}x+\eta,\quad x(0)=\alpha


\end{equation}


имеет единственное решение $x\in{\spc{D}}_0$ при каждой паре


$\{\eta,\alpha\}\in\spc{B}\times\spc{R}^n$ и


это решение непрерывно по норме пространства $\spc{D}_0$


зависит от $\eta$ и $\alpha$ (т.\,е.


для любого $\varepsilon>0$


найдется такое $\delta=\delta(x,\varepsilon)>0$, что


$\|x_1-x\|_{\bold{D}_0}<\varepsilon$, если


$\|\eta_1-\eta\|_{\bold B}<\delta$,


$|\alpha_1-\alpha|<\delta$, где $x_1$ --- решение задачи (3) при


$\eta=\eta_1$, $\alpha=\alpha_1$).


\end{defn}





\begin{thm}


Следующие утверждения эквивалентны$:$


\begin{itemize}


\item[а)] уравнение ${\cal L}x={\cal F}x$\; $\spc{D}_0$-устойчиво$;$


\item[б)] уравнение


\begin{equation}


x={\cal W}({\cal L}_0-{\cal L})x+{\cal W}{\cal F}x+g


\end{equation}


имеет единственное решение $x\in\spc{D}_0$ при каждом


$g\in\spc{D}_0$ и это


решение непрерывно зависит от $g$ по норме пространства $\spc{D}_0;$


\item[в)] уравнение


\begin{equation}


z=({\cal L}_0-{\cal L}){\cal W}z+{\cal F}({\cal W}z+{\cal U}\alpha)


+\vartheta


\end{equation}


имеет единственное решение $z\in\spc{B}$ при каждой паре


$\{\vartheta,\alpha\}\in\spc{B}\times\spc{R}^n$ и это


решение непрерывно зависит от $\vartheta$ и $\alpha$ по норме


пространства $\spc{B}$.


\end{itemize}


\end{thm}





\begin{pf}


Решение $x\in\spc{D}_0$ задачи (3) удовлетворяет


равенству ${\cal L}_0x+{\cal L}x={\cal F}x+\eta+{\cal L}_0x$


и, следовательно, равенству


$x={\cal W}({\cal L}_0-{\cal L})x+{\cal W}{\cal F}x+{\cal W}\eta+


{\cal U}\alpha$. Таким образом, импликация a)\,$\Longrightarrow$\,б)


следует из


того, что при любом $g\in\spc{D}_0$ решение задачи (3), где


$\eta={\cal L}_0g$, $\alpha=g(0)$, является решением уравнения (4).


Импликация б)\,$\Longrightarrow$\,а) следует из того, что при любых


$\eta\in\spc{B}$, $\alpha\in\spc{R}^n$ решение уравнения (4), где


$g={\cal W}\eta+{\cal U}\alpha$, является решением задачи (3).





Импликация а) $\Longrightarrow$ в) есть следствие того, что при любых


$\vartheta\in\spc{B}$, $\alpha\in\spc{R}^n$ решение $z\in\spc{B}$


уравнения (5)


определяется равенством $z={\cal L}_0x$, где $x\in\spc{D}_0$ ---


решение


задачи (3) при $\eta=\vartheta+{\cal L}{\cal U}\alpha$.


Импликация в)\,$\Longrightarrow$\,а) есть следствие того,


что при любых $\eta\in\spc{B}$, $\alpha\in\spc{R}^n$ решение


$x\in\spc{D}_0$


задачи~(3) определяется равенством $x={\cal W}z+{\cal U}\alpha$,


где $z\in\spc{B}$ --- решение уравнения (5) при $\vartheta=


\eta-{\cal L}{\cal U}\alpha$.


\end{pf}





\begin{note}


В предположении


действия операторов $\cal L$ и $\cal F$ из пространства


$\spc{D}_0$ в


пространство $\spc{B}$ можно утверждать, что


$\spc{D}_0$-устойчивость


уравнения ${\cal L}x={\cal F}x$ гарантирует совпадение множества


всех решений уравнения


${\cal L}x={\cal F}x+\eta$ при всех $\eta\in\spc{B}$, которое


обозначим через


$\spc{D}({\cal L}-{\cal F},\spc{B})$,


с множеством $\spc{D}_0$. Действительно, если


$x\in\spc{D}({\cal L}-{\cal F},\spc{B})$, то $x\in\spc{D}_0$ в силу


$\spc{D}_0$-устойчивости. Если же $x\in\spc{D}_0$, то


$x\in\spc{D}({\cal L}-{\cal F},\spc{B})$, т.\,к. элемент $x$ является


решением задачи (3) при $\eta={\cal L}x-{\cal F}x$, $\alpha=x(0)$.


\end{note}





Решения скалярной задачи


$\Dot x(t)+x(t)=x^2(t)$, $x(0)=\alpha$ имеют вид


$x(t)=\alpha/(\alpha-(\alpha-1)e^t)$. При $\alpha<1$ решения этой


задачи


асимптотически устойчивы, но при $\alpha>1$ задача вообще не имеет


решений,


определенных на полуоси $[0,\infty)$.


Поэтому естественно ввести следующее определение ``локальной''


устойчивости. Пусть $x_k$ --- решение задачи (3) при $\eta=\eta_k$,


$\alpha=\alpha_k$, $k\in\spc{Z}_+\eqdef\{0,1,2,\dotsc\}$.





\begin{defn}


Будем говорить, что уравнение ${\cal L}x={\cal F}x+\eta_0$\;


{\it $\spc{D}_0$-устойчиво $($локально$)$ в окрестности


решения} $x_0$ при $x_0(0)=\alpha_0$, если существует


$\delta_0=\delta_0(x_0)>0$, для которого при каждой паре


$\{\eta,\alpha\}\in\spc{B}\times\spc{R}^n$,


$\|\eta-\eta_0\|_{\bold B}\le\delta_0$,


$|\alpha-\alpha_0|\le\delta_0$, задача (3) имеет единственное


решение $x\in\spc{D}_0$, и


это решение непрерывно по норме пространства $\spc{D}_0$


зависит от $\eta$ и $\alpha$


(т.\,е. для любого $\varepsilon>0$


найдется такое $\delta=\delta(x,\varepsilon)>0$, что


$\|x_1-x\|_{\bold{D}_0}<\varepsilon$,


если $\|\eta_1-\eta_0\|_{\bold B}\le\delta_0$,


$|\alpha_1-\alpha_0|\le\delta_0$ и


$\|\eta_1-\eta\|_{\bold B}<\delta$,


$|\alpha_1-\alpha|<\delta$).


\end{defn}





Пусть $\spc{V}$ --- некоторое банахово пространство,


$\Omega\subseteq\spc{V}$ --- область в пространстве $\spc{V}$.


Следуя установившейся в Анализе терминологии,


будем говорить, что уравнение $y={\cal G}y+g$


{\it корректно разрешимо в пространстве} $\spc{V}$


({\it корректно разрешимо в области} $\Omega$),


если оно имеет единственное решение $y\in\spc{V}$ при каждом


$g\in\spc{V}$


($g\in\Omega$) и это решение непрерывно зависит от $g$ по норме


пространства $\spc{V}$.





Отметим, что если оператор $\cal G$ действует в пространстве $\spc{V}$


и уравнение $y={\cal G}y+g$ однозначно разрешимо в этом пространстве


при любом $g\in\spc{V}$, то множество всех решений этого уравнения


при всех $g\in\spc{V}$ совпадает с пространством $\spc{V}$.





\begin{lem}


Пусть уравнение ${\cal L}x=f$ с линейным оператором


${\cal L}:\spc{D}_0\to\spc{B}$\; $\spc{D}_0$--устойчиво. Тогда


\begin{itemize}


\item[а)] $\spc{D}_0$-устойчивость уравнения ${\cal L}x={\cal


F}x+\eta$


эквивалентна корректной разрешимости уравнения $x={\cal C}{\cal


F}x+g$ в пространстве $\spc{D}_0;$


\item[б)] $\spc{D}_0$-устойчивость уравнения ${\cal L}x={\cal


F}x+\eta_0$


в окрестности решения $x_0$


при $\eta_0={\cal L}g_0$, $x_0(0)=\alpha_0\equiv g_0(0)$


эквивалентна корректной разрешимости в пространстве $\spc{D}_0$


уравнения


$x={\cal C}{\cal F}x+g$ для всех $g$ из некоторой окрестности


элемента


$g_0={\cal C}\eta_0+{\cal X}\alpha_0$.


\end{itemize}


\end{lem}





\begin{pf}


$\spc{D}_0$-устойчивость линейного уравнения


${\cal L}x=f$ эквивалентна тому, что его оператор Коши $\cal C$


действует


из пространства $\spc{B}$ в пространство $\spc{D}_0$ и ограничен, а


столбцы


фундаментальной матрицы $X$ принадлежат пространству $\spc{D}_0$.


Поэтому


оператор ${\cal L}:\spc{D}_0\to\spc{B}$ ограничен в силу теоремы


Банаха об


обратном операторе


(\cite{11}, гл. 3, 3.5.3, c.\,89).





Утверждение а) леммы следует из теоремы 1, если положить


${\cal L}_0={\cal L}$.





Докажем утверждение б). Пусть уравнение ${\cal L}x={\cal


F}x+\eta_0$\;


$\spc{D}_0$-устойчиво в некоторой окрестности решения


$x_0$ при $x_0(0)=\alpha_0$.


Это означает существование $\delta_0=\delta_0(x_0)>0$, для которого


при каждой


паре $\{\eta,\alpha\}\in\spc{B}\times\spc{R}^n$,


$\|\eta-\eta_0\|_{\bold B}\le\delta_0$, $|\alpha-\alpha_0|\le\delta_0$,


задача~(3) имеет единственное решение $x\in\spc{D}_0$ и это


решение


непрерывно по норме $\spc{D}_0$ зависит от $\{\eta,\alpha\}$. Пусть


$g,g_0\in\spc{D}_0$, $g={\cal C}\eta+{\cal X}\alpha$,


$g_0={\cal C}\eta_0+{\cal X}\alpha_0$, $\|g-g_0\|_{\bold{D}_0}\equiv


\|\eta-\eta_0\|_{\bold B}+|\alpha-\alpha_0|<\delta_0$. Тогда


$\|\eta-\eta_0\|_{\bold B}\le\delta_0$, $|\alpha-\alpha_0|\le\delta_0$


и при


$\eta={\cal L}g$, $\alpha=g(0)$ задача (3) имеет


единственное решение $x\in\spc{D}_0$, непрерывно зависящее в метрике


$\spc{D}_0$ от


$\{\eta,\alpha\}\in\spc{B}\times\spc{R}^n$. Для решения


$x\in\spc{D}_0$ в


силу уравнения ${\cal L}x={\cal F}x+\eta$ имеем ${\cal


F}x\in\spc{B}$.


Поэтому $x$ удовлетворяет уравнению $x={\cal C}{\cal F}x+g$ и


в метрике $\spc{D}_0$ непрерывно


зависит от $g\in\spc{D}_0$,


$\|g-g_0\|_{\bold{D}_0}\le\delta_0$.





Пусть уравнение $x={\cal C}{\cal F}x+g$ корректно разрешимо в


пространстве $\spc{D}_0$ для $g\in\spc{D}_0$,


$\|g-g_0\|_{\bold{D}_0}\le\delta_0=\delta_0(g_0)>0$.


Положим $\eta_0={\cal L}g_0$, $\alpha=g_0(0)$ и возьмем


$\{\eta,\alpha\}\in\spc{B}\times\spc{R}^n$,


$\|\eta-\eta_0\|_{\bold B}<\delta/2$, $|\alpha-\alpha_0|<\delta/2$.


Тогда


$\|g-g_0\|_{\bold{D}_0}\equiv\|\eta-\eta_0\|_{\bold B}+


|\alpha-\alpha_0|\le\delta_0$,


где $\eta={\cal L}g$, $\alpha=g(0)$, и при таких $g\in\spc{D}_0$


уравнение


$x={\cal C}{\cal F}x+g$ корректно разрешимо в пространстве


$\spc{D}_0$.


Для решения $x\in\spc{D}_0$ в силу уравнения $x={\cal C}{\cal F}x+g$


имеем


${\cal C}{\cal F}x\in\spc{D}_0$, откуда ${\cal F}x\in\spc{B}$.


Поэтому $x$


является решением задачи Коши (3) и в метрике $\spc{D}_0$ 


непрерывно зависит от


$\{\eta,\alpha\}\in\spc{B}\times\spc{R}^n$.


\end{pf}





При исследовании локальной устойчивости удобно пользоваться


подстановкой $y=x-x_0$, сводящей задачу (3) при


$\eta=\eta_0$, $\alpha=\alpha_0$ и $x=x_0$ к ``канонической


форме''


${\cal L}y={\cal F}_0y+\eta_0$, $y(0)=0$, где


${\cal F}_0y\eqdef{\cal F}(y+x_0)-{\cal L}x_0$, и рассматривать


вопрос об


устойчивости в окрестности тривиального решения $y=0$ этой задачи.





Тот факт, что явный вид


оператора ${\cal F}_0$ неизвестен, поскольку неизвестно решение


$x_0$,


не мешает применять здесь стандартные методы: оператор ${\cal F}_0$


наследует


от оператора $\cal F$ требуемые свойства. Например, константы Липшица


операторов $\cal F$ и ${\cal F}_0$ одинаковы.





Некоторые задачи о локальной устойчивости удается решить на основе


следующего


распространения теоремы Ляпунова об устойчивости по первому


приближению, аналогичное теореме~1


из \cite{7} и теоремам 3 из \cite{8}, \cite{9}.





\begin{thm}


Пусть уравнение ${\cal L}x=f$ с линейным


оператором


${\cal L}:{\spc{D}}_0\to{\spc{B}}$\; ${\spc{D}}_0$--устойчиво,


оператор ${\cal F}:\spc{D}_0\to\spc{B}$ обладает свойством$:$ ${\cal


F}(0)=0$ и для любого


$k>0$ найдется такое $\delta>0$, что


\begin{equation}


\|{\cal F}x_2-{\cal F}x_1\|_{\bold B}\le


k\|x_2-x_1\|_{\bold{D}_0}


\end{equation}


при всех $\|x_1\|_{\bold{D}_0}\le\delta$,


$\|x_2\|_{\bold{D}_0}\le\delta$. Тогда


уравнение ${\cal L}x={\cal F}x$\; $\spc{D}_0$-устойчиво


в окрестности тривиального решения.


\end{thm}





\begin{pf}


В силу леммы 1 локальная


$\spc{D}_0$-устойчивость


уравнения ${\cal L}x={\cal F}x$ в окрестности тривиального решения


эквивалентна корректной


разрешимости в пространстве $\spc{D}_0$ уравнения


\begin{equation}


x={\cal C}{\cal F}x+g


\end{equation}


для $g$ из некоторой окрестности нуля. Здесь ${\cal


C}:\spc{B}\to\spc{D}_0$ ---


оператор Коши уравнения ${\cal L}x=f$, $g={\cal C}\eta+{\cal


X}\alpha$.


Обозначим $\sigma\eqdef\|{\cal C}\|_{\bold{B}\to\bold{D}_0}$. В силу


условия на


оператор $\cal F$ для некоторого положительного $k<1/\sigma$ найдется


такое


$\delta>0$, что при любых $x_i\in\spc{D}_0$,


$\|x_i\|_{\bold{D}_0}<\delta$,


$i=1,2$, имеет место неравенство (6).





При $\|g\|_{\bold{D}_0}\le\delta_0\eqdef\delta(1-\sigma k)$ для


уравнения (7)


в шаре $\{\|x\|_{\bold{D}_0}\le\delta_0\}$ выполнены условия


локального принципа


Банаха о


сжимающих отображениях (\cite{11}, гл.\,4, cлед.~4.3.5, c.\,130).


Следовательно, уравнение (4) (которое в данном случае совпадает с


уравнением (7)) корректно разрешимо в пространстве $\spc{D}_0$ при


$\|g\|\le\delta_0$. Отсюда следует $\spc{D}_0$-устойчивость


уравнения ${\cal L}x={\cal F}x$ в окрестности тривиального решения.


\end{pf}





\begin{note}


Условие теоремы относительно оператора $\cal F$


выполнено, если в некоторой окрестности точки $x=0$ оператор


${\cal F}:\spc{D}_0\to\spc{B}$ имеет непрерывную производную Фреше


${\cal F}'$ и ${\cal F}(0)={\cal F}'(0)=0$.


\end{note}





\begin{exam}


Рассмотрим скалярное уравнение


\begin{equation}


\begin{array}{c}


\Dot x(t)+p_0(t)x[h_0(t)]=


p(t)\Dot x^2[h(t)]+q(t)\Dot


x[h(t)]x[h(t)]+r(t)x^2[h(t)]


+f(t),\quad t\ge0,\\[2mm]


x(\xi)=\varphi(\xi),\ \;\Dot x(\xi)=\psi(\xi),\ \text{ если }\ \xi<0.


\end{array}


\end{equation}


Здесь функции $p_0,p,q,r,f:[0,\infty)\to\spc{R}$ измеримы и ограничены


в существенном на $[0,\infty)$; функции $h_0,h:[0,\infty)\to\spc{R}$


измеримы, причем существуют такие $\delta>0$, $b>0$ , что


$t-\delta\le h_0(t)\le t$, $t-\delta\le h(t)\le t$


при всех $t\in[b,\infty)$; $\vraiinf\limits_{t\ge b}p_0(t)>0$,


$\vraisup\limits_{t\ge b}p_0(t)<3/(2\delta)$; начальные функции


$\varphi, \psi:[-\delta,0]\to\spc{R}$ измеримы и ограничены в


существенном


на $[-\delta,0]$.





Следуя \cite{3}--\cite{5}, введем для функций


$y,u:[0,\infty)\to\spc{R}$ и $\phi:\spc{R}\to\spc{R}$


обозначения


$$


y_u(t)\eqdef


\begin{cases}


y[u(t)],&\text{если }u(t)\ge 0;\\0,&\text{если }u(t)<0,


\end{cases}


\qquad


\phi^u(t)\eqdef


\begin{cases}


0,&\text{если }u(t)\ge 0;\\ \phi[u(t)],&\text{если }


u(t)<0.


\end{cases}


$$


Тогда уравнение (8) принимает вид ${\cal L}x={\cal F}x+\eta$, где


\begin{align*}


({\cal L}x)(t)&\eqdef\Dot x(t)+p_0(t)x_{h_0}(t),\\


%


({\cal F}x)(t)&\eqdef p(t)\Dot x^2_h(t)+q(t)\Dot


x_h(t)x_h(t)+r(t)x_h^2(t),\\


%


\eta(t)&\eqdef p(t)\big(\psi^h(t)\big)^2+q(t)\psi^h(t)\varphi^h(t)+


r(t)\big(\varphi^h(t)\big)^2-p_0(t)\varphi^{h_0}(t)+f(t),


\end{align*}


т.\,к. $x[h_0(t)]=x_{h_0}(t)+\varphi^{h_0}(t)$,


$\Dot x[h(t)]=\Dot x_h(t)+\psi^h(t)$,


$x[h(t)]=x_h(t)+\varphi^h(t)$,


$\Dot x^2[h(t)]=\Dot x^2_h(t)+\big(\psi^h(t)\big)^2$,


$\Dot x[h(t)]x[h(t)]=\Dot x_h(t)x_h(t)+\psi^h(t)\varphi^h(t)$,


$x^2[h(t)]=x^2_h(t)+\big(\varphi^h(t)\bigr)^2$.





В силу результатов работ \cite{12}, \cite{13} для функции Коши


$C(t,s)$ уравнения ${\cal L}x=f$ для некоторых $N$,\,$\beta>0$


справедлива при $0\le s\le t$ оценка


$$


|C(t,s)|\le


Ne^{-\beta(t-s)}.


$$


Пусть ${\cal L}_0 x\equiv\Dot x+\beta x=z$ и


$\spc{B}=\spc{L}^{\infty}_\gamma$, $0<\gamma<\beta$.


Тогда банаховы пространства


$\spc{D}_0$, $\spc{D}({\cal L}, \spc{L}^{\infty}_\gamma)$ и


$\spc{W}_\gamma$ совпадают, т.\,е. совпадают как линейные


пространства, и нормы этих пространств эквивалентны, причем


$\spc{W}_\gamma\subset\spc{C}_\gamma$ и это вложение непрерывно.





Здесь через


$\spc{L}^{\infty}_\gamma$, $\gamma\in\spc{R}$, обозначим банахово


пространство всех таких функций $z:[0,\infty)\to\spc{R}^n$, для


которых справедливо представление


$z(t) = y(t)e^{-\gamma t}$, где


$y\in\spc{L}^{\infty}$, с нормой


$\|z\|_{\bold{L}^{\infty}_\gamma}\eqdef\|y\|_{\bold{L}^{\infty}}$,


$\spc{L}_{\infty}$ --- банахово пространство измеримых и ограниченных


в существенном на $[0,\infty)$ функций $y:[0,\infty)\to\spc{R}^n$ с


нормой $\|y\|_{\bold{L}^{\infty}}\eqdef \vraisup\limits_{t\ge


0}|y(t)|$. Аналогично $\spc{W}_\gamma$, $\gamma\in\spc{R}$, ---


банахово пространство всех таких функций $x\in\spc{D}$, для каждой из


которых справедливы включения $x\in\spc{C}_\gamma$, $\Dot


x\in\spc{L}^{\infty}_\gamma$, где


$\|x\|_{\bold{W}_\gamma}\eqdef\|x\|_{\bold{C}_\gamma}+


\|\Dot x\|_{\bold{L}^{\infty}_\gamma}$.





В каждой точке $x_0\in\spc{D}_0$


для любого $x\in\spc{D}_0$ справедливы равенства


$$


\big({\cal F}'(x_0)x\big)(t)=2p(t)(\Dot x_0)_h(t)\Dot x_h(t)+


q(t)[(\Dot x_0)_h(t)x_h(t)+(x_0)_h(t)\Dot


x_h(t)]+2r(t)(x_0)_h(t)x_h(t).


$$


Очевидно, ${\cal F}(0)={\cal F}'(0)=0$.





Таким образом, в силу теоремы 3


уравнение ${\cal L}x={\cal F}x$ локально


$\spc{D}_0$-устойчиво в окрестности тривиального решения.


Из непрерывности вложения


$\spc{D}_0\subset\spc{C}_\gamma$


следует экспоненциальная устойчивость по Ляпунову и, следовательно,


экспоненциальная


устойчивость такого решения по начальным функциям $\varphi$ и


$\psi$, т.\,е. при некотором $M>0$ и достаточно малых


$|\alpha|$ и $\|\varphi\|_{\bold{L}^{\infty}{[-\delta,0]}}\eqdef


\vraisup\limits_{t\in[-\delta,0]}|\varphi(t)|$,


$\|\psi\|_{\bold{L}^{\infty}{[-\delta,0]}}\eqdef


\vraisup\limits_{t\in[-\delta,0]}|\psi(t)|$ справедливо неравенство


$$


|x(t)|\le Me^{-\gamma


t}\big((\|\varphi\|_{\bold{L}^{\infty}{[-\delta,0]}}+


\|\psi\|_{\bold{L}^{\infty}{[-\delta,0]}})^2+|\alpha|\big),\quad


t\ge0.


$$


\end{exam}





Если ограничиться случаем, когда оператор


${\cal F}:\spc{D}_0\to\spc{B}$ допускает непрерывное расширение


${\cal F}:\spc{C}\to\spc{B}$


на пространство $\spc{C}$ непрерывных функций, то чебышевская норма


и удобная естественная упорядоченность ($x_1\le x_2$, если $x_1(t)\le


x_2(t)$ при всех $t\in[0,\infty)$) позволяет получать требуемые


оценки простыми приемами.





Зафиксируем банахово пространство функций


$\spc{B}\subset{\spc{L}}$,


модельное уравнение ${\cal L}_0x=z$ и некоторое банахово пространство


$\spc{V}$ функций $y:[0,\infty)\to{\spc{R}^n}$. Пусть,


далее, $x_k$ --- решение задачи~(3) при


$\eta=\eta_k$, $\alpha=\alpha_k$, $k\in\spc{Z}_+$.





\begin{defn}


Уравнение ${\cal L}x={\cal F}x$ называется


{\it $\spc{V}$-устойчивым}


(обладает {\it $\spc{V}$-свойством}),


если для каждой пары


$\{\eta,\alpha\}\in\spc{B}\times\spc{R}^n$ задача (3) имеет


единственное решение $x\in\spc{V}$ и


это решение непрерывно зависит от $\eta$ и $\alpha$ (т.\,е.


для любого $\varepsilon>0$ найдется


такое $\delta=\delta(x,\varepsilon)>0$, что


$\|x_1-x\|_{\bold V}<\varepsilon$, если


$\|\eta_1-\eta\|_{\bold B}<\delta$,


$|\alpha_1-\alpha|<\delta$).


\end{defn}





\begin{defn}


Будем говорить, что уравнение ${\cal L}x={\cal F}x+\eta_0$\;


{\it $\spc{V}$-устойчиво $($локально$)$ в окрестности решения $x_0$}


при $x_0(0)=\alpha_0$, если существует


$\delta_0=\delta(x_0)>0$, для которого при каждой паре


$\{\eta,\alpha\}\in\spc{B}\times\spc{R}^n$,


$\|\eta-\eta_0\|_{\bold B}\le\delta_0$,


$|\alpha-\alpha_0|\le\delta_0$, задача (3) имеет единственное


решение


$x\in\spc{V}$ и это решение непрерывно зависит от $\eta$ и $\alpha$


(т.\,е. для любого $\varepsilon>0$


найдется такое $\delta=\delta(x,\varepsilon)>0$, что


$\|x_1-x\|_{\bold V}<\varepsilon$, если


$\|\eta_1-\eta_0\|_{\bold B}\le\delta_0$,


$|\alpha_1-\alpha_0|\le\delta_0$ и


$\|\eta_1-\eta\|_{\bold B}<\delta$,


$|\alpha_1-\alpha|<\delta$).


\end{defn}





\begin{note}


Таким образом, например,


$\spc{C}$-устойчивость ($\spc{C}_0$-устойчивость,


$\spc{C}_\gamma$-устой\-чи\-вость), как и локальная


$\spc{C}$-устойчивость ($\spc{C}_0$-устойчивость,


$\spc{C}_\gamma$-устойчивость), гарантируют устойчивость по


Ляпунову либо в целом, либо локально соответственно.


\end{note}





\begin{lem}


Пусть пространствa $\spc{D}_0$ и $\spc{V}$ таковы,


что имеет место непрерывное вложение $\spc{D}_0\subseteq\spc{V}$,


оператор


$\cal F$ непрерывно действует из пространства $\spc{V}$ в


пространство


$\spc{B}$, а уравнение ${\cal L}x=f$ c линейным оператором


${\cal L}:\spc{D}_0\to\spc{B}$\; $\spc{D}_0$-устойчиво. Тогда


\begin{itemize}


\item[а)] $\spc{V}$-устойчивость и $\spc{D}_0$-устойчивость


уравнения ${\cal L}x={\cal F}x$ --- понятия эквивалентные$;$


\item[б)] $\spc{V}$-устойчивость уравнения ${\cal L}x={\cal


F}x+\eta_0$ в окрестности решения $x_0$


и $\spc{D}_0$-устойчивость уравнения ${\cal L}x={\cal


F}x+\eta_0$ в окрестности решения $x_0$


--- понятия


эквивалентные.


\end{itemize}


\end{lem}





\begin{pf}


Из доказательства леммы 1 следует, что операторы


${\cal L}:\spc{D}_0\to\spc{B}$, ${\cal C}:\spc{B}\to\spc{D}_0$,


${\cal X}:\spc{R}^n\to\spc{D}_0$ ограничены. Далее, из непрерывности


вложения


$\spc{D}_0\subseteq\spc{V}$ следует ограниченность оператора Коши


${\cal C}:\spc{B}\to\spc{V}$ и оператора ${\cal


X}:\spc{R}^n\to\spc{V}$.





а) Пусть уравнение ${\cal L}x={\cal F}x$\;


$\spc{D}_0$-устойчиво. Тогда ввиду


вложения $\spc{D}_0\subseteq\spc{V}$ при любых


$\eta=\eta_0\in\spc{B}$ и


$\alpha=\alpha_0\in\spc{R}^n$ решение $x_0\in\spc{D}_0$ задачи (3)


принадлежит


пространству $\spc{V}$, причем другого решения $x\in\spc{V}$ задача (3)


при $\eta=\eta_0$, $\alpha=\alpha_0$ иметь не может. При условии


$\lim\limits_{k\to\infty}\|\eta_k-\eta_0\|_{\bold B}=


\lim\limits_{k\to\infty}|\alpha_k-\alpha_0|=0$ из


$\spc{D}_0$-устойчивости


уравнения ${\cal L}x={\cal F}x$ и непрерывности вложения


$\spc{D}_0\subseteq\spc{V}$ имеем


$\lim\limits_{k\to\infty}\|x_k-x_0\|_{\bold V}\le


d\lim\limits_{k\to\infty}\|x_k-x_0\|_{\bold{D}_0}=0$,


где $d$ --- константа вложения $\spc{D}_0$ в $\spc{V}$.





Пусть уравнение ${\cal L}x={\cal F}x$\; $\spc{V}$-устойчиво.


Тогда при любых


$\eta=\eta_0\in\spc{B}$, $\alpha=\alpha_0\in\spc{R}^n$ решение


$x_0\in\spc{V}$


задачи (3) принадлежит пространству $\spc{D}_0$, т.\,к.


$x_0={\cal C}f_0+{\cal X}\alpha_0$, где $f_0={\cal


F}x_0+\eta_0\in\spc{B}$.


При условии


$\lim\limits_{k\to\infty}\|\eta_k-\eta_0\|_{\bold B}=


\lim\limits_{k\to\infty}|\alpha_k-\alpha_0|=0$ из непрерывности


операторов


${\cal F}:\spc{V}\to\spc{B}$ и ${\cal C}:\spc{B}\to\spc{D}_0$ имеем


$\lim\limits_{k\to\infty}\|x_k-x_0\|_{\bold{D}_0}\le


\|{\cal X}\|_{\bold{R}^n\to{\bold{D}_0}}\lim\limits_{k\to\infty}


|\alpha_k-\alpha_0|+


\|{\cal C}\|_{\bold{B}\to{\bold{D}_0}}


\lim\limits_{k\to\infty}(\|{\cal F}x_k-{\cal F}x_0\|_{\bold B}+


\|\eta_k-\eta_0\|_{\bold B})=0$.





Доказательство утверждения б) производится аналогично, если


рассматривать только такие пары


$\{\eta,\alpha\}\in\spc{B}\times\spc{R}^n$,


которые удовлетворяют неравенствам


$\|\eta-\eta_0\|_{\bold B}\le\delta_0$,


$|\alpha-\alpha_0|\le\delta_0$.


\end{pf}





\begin{note}


В условиях леммы 2\; $\spc{V}$-устойчивость


($\spc{D}_0$-устойчивость) уравнения ${\cal L}x={\cal F}x+\eta$


гарантирует, что множество $\spc{D}({\cal L}-{\cal F},\spc{B})$ всех


решений $x$ этого уравнения при всех $\eta\in\spc{B}$ совпадает с


множеством $\spc{D}_0$. Действительно, при


доказательстве п.\,а) показано, что $\spc{D}({\cal L}-{\cal


F},\spc{B})\subseteq\spc{D}_0$.


Вложение $\spc{D}({\cal L}-{\cal F},\spc{B})\supseteq\spc{D}_0$


очевидно.


\end{note}





\begin{lem}


В условиях леммы $2$


\begin{itemize}


\item[а)] корректная разрешимость уравнения $x={\cal C}{\cal F}x+g$ в


пространстве


$\spc{V}$ гарантирует $\spc{D}_0$-устойчивость уравнения


${\cal L}x={\cal F}x;$


\item[б)] корректная разрешимость уравнения $x={\cal C}{\cal F}x+g$ в


пространстве


$\spc{V}$ при $\|g-g_0\|_{\bold V}\le\delta_0$, где


$g_0\in\spc{D}_0$, $\delta_0>0$,


гарантирует $\spc{D}_0$-устойчивость уравнения


${\cal L}x={\cal F}x+\eta_0$


в некоторой окрестности решения $x_0$ при $x_0(0)=\alpha_0=g_0(0)$,


$\eta_0={\cal L}g_0$.


\end{itemize}


\end{lem}





\begin{pf}


a) Пусть уравнение


$x={\cal C}{\cal F}x+g$ корректно разрешимо в пространстве $\spc{V}$.


Тогда при любом $g\in\spc{D}_0$ решение $x$ этого уравнения


принадлежит


пространству $\spc{D}_0$, т.\,к. $x\in\spc{V}$, ${\cal


F}x\in\spc{B}$ и


${\cal C}{\cal F}x\in\spc{D}_0$.





Возьмем элементы $\eta_0\in\spc{B}$ и $\alpha_0\in\spc{R}^n$, тогда


элемент $g_0={\cal C}\eta_0+{\cal X}\alpha_0\in\spc{D}_0$. Решение


$x_0$


уравнения $x={\cal C}{\cal F}x+g_0$ принадлежит пространству


$\spc{D}_0$,


поэтому $x_0$ является решением задачи ${\cal L}x={\cal F}x+\eta_0$,


$x(0)=\alpha_0$.





Пусть $\{\eta_k,\alpha_k\}\in\spc{B}\times\spc{R}^n$ и


$\lim\limits_{k\to\infty}\|\eta_k-\eta_0\|_{\bold B}=0$,


$\lim\limits_{k\to\infty}|\alpha_k-\alpha_0|=0$,


$g_k={\cal C}\eta_k+{\cal X}\alpha_k$, $k\in\spc{Z}_+$, тогда


$\lim\limits_{k\to\infty}\|g_k-g_0\|_{\bold{D}_0}=0$ ввиду


$\spc{D}_0$-устойчивости линейного уравнения ${\cal L}x=f$.





Из непрерывности


вложения $\spc{D}_0\subseteq\spc{V}$ и корректной разрешимости


уравнения


$x={\cal C}{\cal F}x+g$ в пространстве $\spc{V}$ получаем


$\lim\limits_{k\to\infty}\|g_k-g_0\|_{\bold V}=0$ и


$\lim\limits_{k\to\infty}\|x_k-x_0\|_{\bold V}=0$. Далее из


непрерывности


операторов ${\cal F}:\spc{V}\to\spc{B}$ и ${\cal


C}:\spc{B}\to\spc{D}_0$


имеем


$\lim\limits_{k\to\infty}\|x_k-x_0\|_{\bold{D}_0}\le


\|{\cal C}\|_{\bold{B}\to{\bold{D}_0}}


\lim\limits_{k\to\infty}\|{\cal F}x_k-{\cal F}x_0\|_{\bold B}+


\lim\limits_{k\to\infty}\|g_k-g_0\|_{\bold{D}_0}=0$.





б) Пусть уравнение $x={\cal C}{\cal F}x+g$ корректно разрешимо в


пространстве


$\spc{V}$ при $g\in\spc{V}$, $\|g-g_0\|_{\bold V}\le\delta_0$, где


$g_0\in\spc{D}_0$, $\delta_0>0$. Как показано в п.\,а), при всех


таких


$g\in\spc{D}_0$ решение $x$ уравнения $x={\cal C}{\cal F}x+g$


принадлежит пространству


$\spc{D}_0$. Возьмем элементы $\eta\in\spc{B}$, $\alpha\in\spc{R}^n$


так, чтобы выполнялись неравенства


$\|\eta-\eta_0\|_{\bold B}\le\delta_0/(2d)$,


$|\alpha-\alpha_0|\le\delta_0/(2d)$, где $d$ --- константа


вложения $\spc{D}_0$


в $\spc{V}$. Тогда для элементов $g={\cal C}\eta+{\cal X}\alpha$


имеем


$\|g-g_0\|_{\bold V}\le


d\|g-g_0\|_{\bold{D}_0}=d(\|\eta-\eta_0\|_{\bold B}+


|\alpha-\alpha_0|)\le\delta_0$. Для таких элементов $g$ уравнение


$x={\cal C}{\cal F}x+g$ корректно разрешимо в пространстве $\spc{V}$


и


каждое его решение $x$ является решением задачи Коши (3).


Непрерывную зависимость по норме $\spc{D}_0$ решения


$x$


этой задачи можно показать аналогично тому, как это сделано в п.\,а).


\end{pf}





\begin{note}


Если выполнены условия леммы 2, то в случае корректной


разрешимости уравнения $x={\cal C}{\cal F}x+g$ в пространстве


$\spc{V}$


множество $\frak M$ всех решений этого уравнения при всех


$g\in\spc{D}_0$ совпадает с множеством


$\spc{D}_0$. Действительно, в п.\,а) доказательства леммы 2


показано, что $\frak M\subseteq\spc{D}_0$. Вложение


$\frak M\supseteq\spc{D}_0$ очевидно.


\end{note}





Справедливо следующее обобщение теоремы 2 из \cite{7} и теоремы 5


из \cite{9}.





\begin{thm}


Пусть пространство $\spc{D}_0$ непрерывно


вложено в пространство $\spc{V}$, уравнение ${\cal L}x=f$ с линейным


оператором ${\cal L}:\spc{D}_0\to\spc{B}$\; $\spc{D}_0$-устойчиво, а


оператор


${\cal F}:\spc{V}\to\spc{B}$ обладает свойством$:$


${\cal F}(0)=0$ и для любого


$k>0$ найдется такое $\delta>0$, что


\begin{equation}


\|{\cal F}x_2-{\cal F}x_1\|_{\bold B}\le


k\|x_2-x_1\|_{\bold V}


\end{equation}


при всех


$\|x_1\|_{\bold V}\le\delta$, $\|x_2\|_{\bold V}\le\delta$. Тогда


уравнение ${\cal L}x={\cal F}x$ локально


$\spc{D}_0$-устойчиво в окрестности тривиального решения.


\end{thm}





\begin{pf}


Из доказательства теоремы 3 следует, что


операторы ${\cal L}:\spc{D}_0\to\spc{B}$, ${\cal


C}:\spc{B}\to\spc{V}$, ${\cal X}:\spc{R}^n\to\spc{V}$ ограничены.


Повторяя схему доказательства теоремы 3, получим, что уравнение


$x={\cal C}{\cal F}x+g$ корректно разрешимо в пространстве $\spc{V}$


при $g\in\spc{V}$, $\|g\|_{\bold V}\le\delta_0$ для некоторого


$\delta_0>0$. Тогда в силу утверждения~б) леммы~3 уравнение


${\cal L}x={\cal F}x$ локально $\spc{D}_0$-устойчиво в


окрестности тривиального решения.


\end{pf}





\begin{lem}


Пусть линейное уравнение ${\cal L}x=f$\;


$\spc{V}$-устойчиво. Тогда


\begin{itemize}


\item[а)] корректная разрешимость 


уравнения $x={\cal C}{\cal F}x+g$ в пространстве


$\spc{V}$ гарантирует $\spc{V}$-устойчивость уравнения


${\cal L}x={\cal F}x;$


\item[б)] корректная разрешимость 


уравнения $x={\cal C}{\cal F}x+g$ в пространстве


$\spc{V}$ при $\|g-g_0\|_{\bold V}\le\delta_0$, где


$g_0={\cal C}\eta_0+{\cal X}\alpha_0$, $\eta_0\in\spc{B}$,


$\alpha_0\in\spc{R}^n$, $\delta_0>0$,


гарантирует $\spc{V}$-устойчивость уравнения


${\cal L}x={\cal F}x+\eta_0$ в некоторой окрестности решения $x_0$ при


$x_0(0)=\alpha_0$.


\end{itemize}


\end{lem}





\begin{pf}


$\spc{V}$-устойчивость уравнения ${\cal


L}x=f$ означает, что оператор Коши $\cal C$ действует из пространства


$\spc{B}$ в пространство $\spc{V}$ и ограничен, а столбцы


фундаментальной матрицы $X$ уравнения принадлежат пространству


$\spc{V}$. Обозначим $\varsigma=\|{\cal C}\|_{\bold{B}\to\bold{V}}+


\|{\cal X}\|_{\bold{R}^n\to\bold{V}}$.


\smallskip





а) Пусть уравнение $x={\cal C}{\cal F}x+g$ корректно разрешимо в


пространстве $\spc{V}$. Рассмотрим


последовательность задач


${\cal L}x={\cal F}x+\eta_i$, $x(0)=\alpha_i$ при


$\{\eta_i,\alpha_i\}\in\spc{B}\times\spc{R}^n$, $i\in\spc{Z}_+$. Тогда


$g_i={\cal C}\eta_i+{\cal X}\alpha_i\in\spc{V}$ и $\lim\limits_{i\to\infty}


\|g_i-g_0\|_{\bold V}=0$, если


$\lim\limits_{i\to\infty}\|\eta_i-\eta_0\|_{\bold B}=


\lim\limits_{i\to\infty}|\alpha_i-\alpha_0|=0$.


Решение $x_i$ уравнения $x={\cal C}{\cal F}x+g_i$ существует, единственно,


причем $x_i\in\spc{V}$ и $\lim\limits_{i\to\infty}\|x_i-x_0\|_{\bold V}=0$


при $\lim\limits_{i\to\infty}\|g_i-g_0\|_{\bold V}=0$. Элементы $x_i$ и $g_i$


связаны равенством $x_i={\cal C}({\cal F}x_i+\eta_i)+{\cal X}\alpha_i$,


поэтому элемент $x_i$ является решением задачи ${\cal L}x=f_i$,


$x(0)=\alpha_i$ при $f_i={\cal F}x_i+\eta_i$. Более того,


$\lim\limits_{i\to\infty}\|x_i-x_0\|_{\bold V}=0$, если


$\lim\limits_{i\to\infty}\|\eta_i-\eta_0\|_{\bold B}=


\lim\limits_{i\to\infty}|\alpha_i-\alpha_0|=0$.


\smallskip





б) Пусть уравнение $x={\cal C}{\cal F}x+g$ корректно разрешимо в


пространстве $\spc{V}$ при $g\in\spc{V}$, $\|g-g_0\|_{\bold V}\le\delta_0$,


где $\delta_0>0$ и $g_0={\cal C}\eta_0+{\cal X}\alpha_0$,


$\{\eta_0,\alpha_0\}\in\spc{B}\times\spc{R}^n$. Рассмотрим


последовательность задач


${\cal L}x={\cal F}x+\eta_i$, $x(0)=\alpha_i$


при


$\{\eta_i,\alpha_i\}\in\spc{B}\times\spc{R}^n$,


$\|\eta_i-\eta_0\|_{\bold B}\le\varepsilon/\varsigma$,


$|\alpha_i-\alpha_0|\le\varepsilon/\varsigma$, $i\in\spc{Z}_+$.


Тогда


$g_i={\cal C}\eta_i+{\cal X}\alpha_i\in\spc{V}$,


$\|g_i-g_0\|_{\bold V}\le\delta_0$. Кроме того,


$\lim\limits_{i\to\infty}\|g_i-g_1\|_{\bold V}=0$, если


$\lim\limits_{i\to\infty}\|\eta_i-\eta_1\|_{\bold B}=


\lim\limits_{i\to\infty}|\alpha_i-\alpha_1|=0$. Повторяя рассуждения


п.\,а), получаем, что решение $x_i$ уравнения $x={\cal C}{\cal F}x+g_i$


является решением задачи ${\cal L}x={\cal F}x+\eta_i$, $x(0)=\alpha_i$,


причем $\lim\limits_{i\to\infty}\|x_i-x_1\|_{\bold V}=0$,


если $\lim\limits_{i\to\infty}\|\eta_i-\eta_1\|_{\bold B}=


\lim\limits_{i\to\infty}|\alpha_i-\alpha_1|=0$.


\end{pf}





Имеет место другой вариант обобщения теоремы 2 из \cite{7} и


теоремы 6 из \cite{9}.





\begin{thm}


Пусть линейное уравнение ${\cal L}x=f$\;


$\spc{V}$-устойчиво, а оператор ${\cal F}:\spc{V}\to\spc{B}$ обладает


свойством$:$ ${\cal F}(0)=0$ и для любого $k>0$ найдется такое $\delta>0$,


что при $\|x_i\|_{\bold V}\le\delta$, $i=1,2$, имеет место неравенство


$(9)$. Тогда уравнение ${\cal L}x={\cal F}x$ локально


$\spc{V}$-устойчиво в некоторой окрестности тривиального


решения.


\end{thm}





\begin{pf}


Повторяя схему доказательства теоремы 3,


установим, что уравнение $x={\cal C}{\cal F}x+g$ корректно разрешимо в


пространстве $\spc{V}$ при $g\in\spc{V}$, $\|g\|_{\bold V}\le\delta_0$


для некоторого $\delta_0>0$. Тогда в силу утверждения~а) леммы~4


уравнение ${\cal L}x={\cal F}x$ локально $\spc{V}$-устойчиво


в некоторой окрестности тривиального решения.


\end{pf}





Для установления локальной $\spc{D}_0$-устойчивости в окрестности решения


уравнения


${\cal L}x={\cal F}x$ на основе теоремы 3 следует разложить оператор


${\cal F}:\spc{D}_0\to\spc{B}$ на сумму вида ${\cal F}={\cal L}_1+{\cal F}_1$,


где ${\cal F}_1$ удовлетворяет условиям теоремы. Если при этом уравнение


$({\cal L}-{\cal L}_1)x=f$ окажется $\spc{D}_0$-устойчивым, то локальная


$\spc{D}_0$-устойчивость в окрестности решения нелинейного уравнения будет


установлена.





Обозначим ${\cal L}_1x\eqdef{\cal F}'(x_0)x$,


${\cal F}_1x={\cal F}x-{\cal L}_1x$.


Тогда уравнение ${\cal L}x={\cal F}x$ можно записать в виде


$$


({\cal L}-{\cal L}_1)x={\cal F}_1x


$$


и в качестве линейного приближения можно


брать уравнение $({\cal L}-{\cal L}_1)x=f$.





Рассмотрим простые примеры применения этой схемы.





\begin{exam}


В скалярном уравнении


\begin{equation}


\begin{array}{c}


\Dot x(t)=f(x[t-\vartheta],\Dot x[t-\tau]),\quad t\ge0,\\[2mm]


x(\xi)=\Dot x(\xi)=0\; \text{ при }\; \xi<0,


\end{array}


\end{equation}


$\vartheta$, $\tau$ --- положительные постоянные, функция $f$ непрерывно


дифференцируема в некоторой окрестности нуля пространства $\spc{R}^2$,


причем $f(0,0)\equiv0$. Обозначим $h(t)=t-\vartheta$, $g(t)=t-\tau$,


${\cal L}x=\Dot x$, ${\cal F}x={\cal N}_f(x_h,\Dot x_g)$, $p=-f_u(0,0)$,


$q=f_v(0,0)$, где ${\cal N}_f(y,z)\equiv f(y,z)$ --- оператор


Немыцкого, порожденный функцией $f(\cdot,\cdot)$,


$f_u(u,v)\equiv\frac{\partial f(u,v)}{\partial u}$,


$f_v(u,v)\equiv\frac{\partial f(u,v)}{\partial


v}$. Выберем модельное уравнение ${\cal L}_0x\equiv\Dot x+x=z$ и


пространство $\spc{B}=\spc{L}^{\infty}$. В указанных предположениях


оператор ${\cal N}_f$ действует из некоторого шара с центром в нуле


пространства $\spc{L}^{\infty}\times\spc{L}^{\infty}$ в пространство


$\spc{L}^{\infty}$ и непрерывно дифференцируем в этом шаре


(\cite{14}, гл.\,X, \S\,1, 1.10, cc.\,385, 386).


Тогда оператор $\cal F$ действует из


некоторого шара с центром в нуле пространства $\spc{D}_0$ в


пространство $\spc{L}^{\infty}$ и непрерывно дифференцируем в этом


шаре, причем ${\cal L}_1x\equiv{\cal F}'(0)x=-px_h+q\Dot x_g$.


Уравнение (10) можно записать в виде $({\cal L}-{\cal L}_1)x={\cal


F}_1x$, где $({\cal L}-{\cal L}_1)x=\Dot x-q\Dot x_g+px_h$, ${\cal


F}_1x={\cal F}x-{\cal L}_1x$.





В силу признаков из \cite{15}, \cite{16} достаточным


условием $\spc{C}$-устойчивости линейного уравнения


$({\cal L}-{\cal L}_1)x=\eta$ является выполнение (при $\omega=e^{-1}$,


$\sigma(\omega)=1$) неравенств $|q|<1$, $p>0$ и


$$


(1-q)|p\vartheta+(q-1)/e|+p|q|\tau<1-3|q|+q^2.


$$





Тогда в этих предположениях из теоремы 3 следует локальная


$\spc{D}_0$-устойчивость уравнения


$({\cal L}-{\cal L}_1)x={\cal F}_1x$ (и уравнения (10))


в окрестности тривиального решения.


\end{exam}





\begin{exam}


В системе уравнений модели


``хищник--жертва'', учитывающей внутривидовую борьбу в популяциях


\cite{17},


\begin{align}


\Dot N_1(t)&=\bigg(\varepsilon_1+\int_{t-\tau_1}^tN_2(s)d_sK_1(t-s)-


\beta_1N_1(t)\bigg)N_1(t),\notag\\


%


\Dot N_2(t)&=\bigg(-\varepsilon_2+\int_{t-\tau_2}^tN_1(s)d_sK_2(t-s)-


\beta_2N_2(t)\bigg)N_2(t),\quad t\ge0,\\


%


N_1(\xi)&=\wt N_1(\xi),\quad N_2(\xi)=\wt N_2(\xi),


\;\text{ если }\; \xi<0,\notag


\end{align}


$N_1$ и $N_2$ --- численности жертв и хищников соответственно;


константы


$\varepsilon_1$, $\varepsilon_2$, $\tau_1$, $\tau_2$, $\beta_1$, $\beta_2$


положительны;


$\wt N_1:[-\tau_2,0]\to[0,\infty)$, $\wt N_2:[-\tau_1,0]\to[0,\infty)$


---


кусочно--непрерывные функции; $K_1:[0,\tau_1]\to\spc{R}$,


$K_2:[0,\tau_2]\to\spc{R}$ --- неубывающие ограниченные функции. Обозначим


$k_1=K_1(\tau_1)-K_1(0)$, $k_2=K_2(\tau_2)-K_2(0)$.





Найдем условия устойчивости нетривиального положения равновесия


$$


N_1^0=\frac{\varepsilon_1\beta_2+\varepsilon_2k_1}{\beta_1\beta_2+k_1k_2},


\quad


N_2^0=\frac{\varepsilon_1k_2-\varepsilon_2\beta_1}{\beta_1\beta_2+k_1k_2}


$$


при $\wt N_1(\xi)\equiv N_1^0$, $\wt N_2(\xi)\equiv N_2^0$.


Заметим, что $N_2^0>0$ при


$\varepsilon_1k_2>\varepsilon_2\beta_1$. Будем предполагать выполненным это


неравенство. Введя новые функции $x_1(t)=N_1(t)-N_1^0\,$,\,


$x_2(t)=N_2(t)-N_2^0$, $\varphi_1(\xi)=\wt N_1(\xi)-N_1^0$,


$\varphi_2(\xi)=\wt N_2(\xi)-N_2^0$, получим систему


\begin{equation}


\begin{split}


\Dot x_1(t)&=(-\gamma_1+a_1(t))x_1(t)+


\int_{h_1(t)}^t x_2(s)d_sr_1(t-s)+


{\cal F}_1(x_1,x_2)(t)+\eta_1(t),\\


%


\Dot x_2(t)&=(-\gamma_2+a_2(t))x_2(t)-


\int_{h_2(t)}^t x_1(s)d_sr_2(t-s)+


{\cal F}_2(x_1,x_2)(t)+\eta_2(t),\quad t\ge0.


\end{split}


\end{equation}


Здесь $\gamma_1=\beta_1N_1^0$, $\gamma_2=\beta_2N_2^0$;


$a_1(t)=\int\limits_{h_1^-(t)}^0\varphi_2(s)d_sK_1(t-s)$,


$a_2(t)=\int\limits_{h_2^-(t)}^0\varphi_1(s)d_sK_2(t-s)$;


$h_1^-(t)=\min\{t-\tau_1,0\}$, $h_2^-(t)=\min\{t-\tau_2,0\}$,


$h_1(t)=\max\{t-\tau_1,0\}$, $h_2(t)=\max\{t-\tau_2,0\}$;


$r_1(t)=N_1^0K_1(t)$, $r_2(t)=N_2^0K_2(t)$, $p_1=N_1^0k_1$,


$p_2=N_2^0k_2$;


\begin{align*}


{\cal F}_1(x_1,x_2)(t)&=


x_1(t)\bigg(\int_{h_1(t)}^tx_2(s)d_sK_1(t-s)-


\beta_1x_1(t)\bigg), \\


%


{\cal F}_2(x_1,x_2)(t)&=


-x_2(t)\bigg(\int_{h_2(t)}^tx_1(s)d_sK_2(t-s)+


\beta_2x_2(t)\bigg);


\end{align*}


$\eta_1(t)=\int\limits_{h_1^-(t)}^0\varphi_2(s)d_sr_1(t-s)$,


$\eta_2(t)=-\int\limits_{h_2^-(t)}^0\varphi_1(s)d_sr_2(t-s)$.


Обозначим $x=\col\{x_1,x_2\}$,


$\eta=\col\{\eta_1,\eta_2\}$,


${\cal F}x=\col\{{\cal F}_1(x_1,x_2),


{\cal F}_2(x_1,x_2)\}$;


$$


({\cal L}_0x)(t)=


\begin{pmatrix}\Dot x_1(t)\\ \Dot x_2(t)\end{pmatrix}+


\begin{pmatrix}\gamma_1 & 0 \\ 0 & \gamma_2\end{pmatrix}


\begin{pmatrix}x_1(t) \\ x_2(t)\end{pmatrix},


$$


${\cal L}={\cal L}_0-{\cal R}$,


${\cal R}x=\col\{{\cal R}_1(x_1,x_2),{\cal R}_2(x_1,x_2)\}$,


где


\begin{align*}


{\cal R}_1(x_1,x_2)(t)&=a_1(t)x_1(t)+


\int_{h_1(t)}^tx_2(s)d_sr_1(t-s),\\


%


{\cal R}_2(x_1,x_2)(t)&=-a_2(t)x_2(t)-


\int_{h_2(t)}^tx_1(s)d_sr_2(t-s).


\end{align*}





Тогда систему (12) можно записать в виде матричного уравнения


$({\cal L}x)(t)=({\cal F}x)(t)+\eta(t)$, $t\ge0$, где оператор


$\cal F$ непрерывно действует из пространства $\spc{C}$ непрерывных и


ограниченных вектор-функций $x:[0,\infty)\to\spc{R}^2$ в пространство


$\spc{L}^{\infty}$ измеримых и ограниченных в существенном


вектор-функций $z:[0,\infty)\to\spc{R}^2$. Более того,


оператор ${\cal F}:\spc{C}\to\spc{L}^{\infty}$ непрерывно


дифференцируем по Фреше в любой точке $x\in\spc{C}$ и ${\cal


F}(0)={\cal F}'(0)=0$. Укажем условия $\spc{C}$-устойчивости


линейного уравнения ${\cal L}x=f$. Тогда в силу теоремы 4


нелинейное уравнение ${\cal L}x = {\cal F}x$ будет также локально


$\spc{C}$-устойчивым в некоторой окрестности тривиального


решения.





Оператор Коши $\cal W$ модельного уравнения ${\cal L}_0x=z$ определен


равенством


$$


({\cal W}z)(t)=\int_0^tW(t,s)z(s)\,ds,


$$


где


$$


W(t,s)=\begin{pmatrix}


e^{-\gamma_1(t-s)} & 0 \\ 0 & e^{-\gamma_2(t-s)}\end{pmatrix}.


$$


Очевидно, оператор $\cal W$ действует из пространства $\spc{L}^{\infty}$ в


пространство $\spc{C}$ и ограничен. Поэтому банахово пространство


$\spc{D}_0\eqdef\spc{D}({\cal L}_0,\spc{L}^{\infty})$ совпадает с


банаховым пространством $\spc{W}$. В частности, норма


$\|x\|_{\bold{D}_0}\eqdef\|{\cal L}_0x\|_{\bold{L}^{\infty}}+|x(0)|$ в


пространстве $\spc{D}_0$


эквивалентна норме $\|x\|_{\bold W}\eqdef\|x\|_{\bold C}+\|\Dot


x\|_{\bold{L}^{\infty}}$ в пространстве $\spc{W}$. Оператор ${\cal R}$


(${\cal G}\eqdef{\cal W}{\cal R}$) действует из пространства


$\spc{C}$ в пространство $\spc{L}^{\infty}$ (в пространство


$\spc{C}$) и ограничен. В силу замечания 4 и теоремы 4 


\cite{1} однозначная разрешимость уравнения $x={\cal G}^bx+g$ в


пространстве $\spc{C}^b\eqdef \{x\in\spc{C}:x(t)\equiv 0$ на


$[0,b]\}$ для некоторого $b>0$ гарантирует


$\spc{C}$-устойчивость ($\spc{D}_0$-устойчивость)


уравнения ${\cal L}x=f$.





Определим норму $|\alpha|$ элемента


$\alpha=\col\{\alpha_1,\alpha_2\}\in\spc{R}^2$


равенством


$$


|\alpha|=\max\{|\alpha_1|,|\alpha_2|\},


$$


норму $\|x\|_{\bold{C}^b}$ элемента


$x\in\spc{C}^b$ --- равенством


$$


\|x\|_{\bold{C}^b}=


\Big|\col\Big\{\sup\limits_{t\ge b}|x_1(t)|,


\,\sup\limits_{t\ge b}|x_2(t)|\Big\}\Big|.


$$


Тогда


$\|{\cal G}\|_{\bold{C}^b\to\bold{C}^b}<1$, если $b>\max\{\tau_1,\tau_2\}$,


$p_1<\gamma_1$, $p_2<\gamma_2$.





Итак, из теоремы 4 (теоремы 3) следует, что система (12) локально


$\spc{C}$-устойчива ($\spc{D}_0$-устойчива) в некоторой


окрестности тривиального решения, если $k_1<\beta_1$, $k_2<\beta_2$.





Покажем, что в условиях $k_1<\beta_1$, $k_2<\beta_2$ система (11)


локально $\spc{C}_\gamma$-устойчива в окрестности положения


равновесия $N_1=N_1^0$, $N_2=N_2^0$ для некоторого $\gamma>0$.





Действительно, пусть $\gamma_1>p_1$, $\gamma_2>p_2$, тогда найдется такое


положительное число $\gamma$, что


$\gamma_1-\gamma>p_1$, $\gamma_2-\gamma>p_2$. Оператор $\cal R$ действует из


пространства $\spc{C}_\gamma$ в пространство


$\spc{L}^{\infty}_\gamma$, т.\,к. для $x\in\spc{C}_\gamma$,


$x(t)=e^{-\gamma t}y(t)$, $y\in\spc{C}$, $y(t)\in \spc{R}^2$, имеем


$\|{\cal R}x\|_{\bold{L}^{\infty}_\gamma}\le


L\|x\|_{\bold{C}_\gamma}$, где


$$


L=\max\Big\{\vraisup_{t\ge0}|a_1(t)|+p_1e^{\gamma\tau_1},\,


\vraisup_{t\ge0}|a_2(t)|+p_2e^{\gamma\tau_2}\Big\}.


$$


Аналогично можно показать действие и непрерывность операторов


${\cal W}:\spc{L}^{\infty}_\gamma\to\spc{C}_\gamma$,


${\cal F}:\spc{C}_\gamma\to\spc{L}^{\infty}_\gamma$.


Условие однозначной разрешимости уравнения $x={\cal W}{\cal R}x+g$ в


пространстве $\spc{C}_\gamma$ имеет вид


$p_1<\gamma_1-\gamma$, $p_2<\gamma_2-\gamma$. Тогда из теоремы 4


(теоремы 3) следует, что система (12) локально


$\spc{C}_\gamma$-устойчива ($\spc{D}({\cal


L}_0, \spc{L}^{\infty}_\gamma)$-устойчива) в окрестности


тривиального решения. Поэтому система (11) локально


$\spc{C}_\gamma$-устойчива ($\spc{D}({\cal L}_0,


\spc{L}^{\infty}_\gamma)$-устойчива) в окрестности положения


равновесия $N_1=N_1^0$, $N_2=N_2^0$ при $k_1<\beta_1$, $k_2<\beta_2$


и $0<\gamma<\min\{(\beta_1-k_1)/N_1^0, (\beta_2-k_2)/N_2^0\}$.


\end{exam}





\begin{note}


В статье \cite{17} условие асимптотической устойчивости по начальным


функциям нетривиального положения равновесия $\{N_1^0,N_2^0\}$


уравнения (11) имеет вид


$N_1^0k_1+N_2^0k_2<2\min\{N_1^0\beta_1,N_2^0\beta_2\}$. Это условие


совпадает с полученными нами условиями в отдельных случаях. В общем


случае условия пересекаются, но не совпадают.


\end{note}








\begin{thebibliography}{99}





\bibitem{1}


Азбелев Н.В., Симонов П.М. {\em Устойчивость уравнений с


запаздывающим аргументом} // Изв. вузов. Математика. -- 1997. --


\No\,6. -- С.\,3--16.


\bibitem{2}


Azbelev N.V., Rakhmatullina L.F.


{\em Stability of solutions of the equations with aftereffect} //


Funct. Different. Equat. -- 1998. -- V.\,5. -- \No\,1--2. --


P.\,39--55.


\bibitem{3}


Азбелев Н.В., Максимов В.П., Рахматуллина Л.Ф. {\em Введение


в теорию функционально--дифференциальных


уравнений.} -- М.: Наука, 1991. -- 280~с.


\bibitem{4}


Azbelev N.V., Maksimov V.P., Rakhmatullina L.F. {\em Introduction


to the theory of linear functional differential equations.} --


Atlanta: World Federation Publ. Company, 1995. -- 172~p.


\bibitem{5}


Azbelev N.V., Rakhmatullina L.F. {\em Theory of linear


abstract functional differential equations and applications} //


Mem. Different. Equat. and Mathem. Physics. -- Tbilisi:


Publ. House GCI, 1996. -- V.\,8. -- P.\,1--102.


\bibitem{6}


Азбелев Н.В., Малыгина В.В. {\em Об устойчивости


тривиального решения нелинейных уравнений с


последействием} // Изв. вузов. Математика. -- 1994. -- \No\,6. --


С.\,72--85.


\bibitem{7}


Азбелев Н.В., Ермолаев М.Б., Симонов П.М. {\em К


вопросу об устойчивости


функционально-дифференциальных


уравнений по первому приближению} // Изв. вузов. Математика. -- 1995.


\No\,10. -- С.\,3--9.


\bibitem{8}


Азбелев Н.В., Симонов П.М. {\em Исследование устойчивости нелинейных


дифференциальных уравнений с


последействием по линейному приближению} // Некоторые пробл.


фундам. и прикл. матем. -- М.: МФТИ(ГУ), 1998.


-- С.\,4--14.


\bibitem{9}


Азбелев Н.В., Симонов П.М. {\em Устойчивость и асимптотическое


поведение решений нелинейных дифференциальных уравнений с


запаздывающим аргументом} // Междунар. конгресс


``Нелинейный анализ и его приложения'': Тр. конгресса. --


М.: ЦИУНД при ИМАШ РАН. -- 1999. -- С.\,658--673.


\bibitem{10}


Красовский Н.Н.


{\em Некоторые задачи теории устойчивости движения.} -- М.: Физматгиз,


1959. -- 212~с.


\bibitem{11}


Хатсон В.К.Л., Пим Дж.С.


{\em Приложения функционального анализа и теории операторов.} -- М.:


Мир, 1983. -- 432~с.


\bibitem{12}


Малыгина В.В. {\em Некоторые признаки устойчивости


уравнений с запаздывающим аргументом} // Дифференц. уравнения. --


1992. -- Т.\,28. -- \No\,10. -- С.\,1716--1723.


\bibitem{13}


Малыгина В.В. {\em Об устойчивости


решений некоторых линейных дифференциальных уравнений с


последействием} // Изв. вузов. Математика. -- 1993. -- \No\,5. --


С.\,72--85.


\bibitem{14}


Забрейко П.П., Кошелев А.И., Красносельский М.А. и др.


{\em Интегральные уравнения.} -- М.: Наука, 1968. -- 448~с.


\bibitem{15}


Гусаренко С.А. {\em Признаки разрешимости задач о


накоплении возмущений для


функ\-ци\-о\-наль\-но--дифференциальных уравнений} //


Функц.--дифференц. уравнения. -- Пермь: Пермск. политехн. ин-т.,


1987. -- С.\,30--40.


\bibitem{16}


Гусаренко С.А. {\em Об ограниченности оператора Коши} //


Функц.--дифференц. уравнения. -- Пермь:


Пермск. политехн. ин-т. -- 1992. -- С.\,111--122.


\bibitem{17}


Дроздов А.Д., Колмановский В.Б., Триджанте Д. {\em Об


устойчивости системы хищник--жертва} // Автоматика и телемеханика. --


1992. -- \No\,11. -- С.\,57--64.





\end{thebibliography}





\vskip 2cm





\post{\it Пермский государственный} {\it Поступила}


\post{\it технический университет,}{24.05.1999}


\smallskip


\post{\it Пермский государственный университет}{}





\label{end}


\end{document}


